% !TeX root = ../main.tex
% ---------------------------------------------------------------------------
% Chapter imported from Weekly Study Notes/ Chapter~\ref{ch:context}.tex
% Initial enrichment pass: copied directly from the weekly study-note body.
% ---------------------------------------------------------------------------
\chapter{Infrastructure and Geotechnical Challenges}
\label{ch:context}

\section{Infrastructure demand and delivery capacity}
\label{sec:infrastructure-demand-and-delivery-capacity}

\subsection{Infrastructure Market Capacity}

\paragraph{Historic Level of Investment}
Australia has experienced a sustained period of major infrastructure investment, including a large number of underground transport, water, and energy projects. This creates strong demand for geotechnical engineers, but also places pressure on investigation programs, design timelines, construction resources, and risk management processes.

\begin{itemize}
 \item The 2021 infrastructure pipeline included \textbf{125 major transport infrastructure projects}, representing approximately \textbf{58\% of the infrastructure pipeline}.
 \item Tunnelling projects were expected to increase substantially over the following five years.
 \item Engineers should compare the 2021 Infrastructure Market Capacity Report with the most recent available report and note changes in project numbers, investment scale, and delivery risks.
\end{itemize}

\begin{notebox}[Teaching note]
The numerical values above provide context rather than design input. They are included to show why geotechnical investigation and risk management are increasingly important.
\end{notebox}

\paragraph{Selected Major Tunnelling Projects in Australia}

\begin{table}[H]
\centering
\small
\begin{tabularx}{\textwidth}{l l X}
\toprule
\textbf{Project} & \textbf{Location} & \textbf{Key features} \\
\midrule
Melbourne Metro Tunnel & Victoria & Approximately 9 km twin tunnels and five new underground stations. \\
Snowy 2.0 & New South Wales & Approximately 27 km waterway tunnel and an underground power station at large depth. \\
West Gate Tunnel & Victoria & Large-diameter road tunnel constructed in an urban environment. \\
Sydney Metro West & New South Wales & Twin metro tunnels constructed using TBMs. \\
Suburban Rail Loop & Victoria & Long-term orbital rail project with major underground sections. \\
North East Link & Victoria & Major road tunnel project in Melbourne. \\
\bottomrule
\end{tabularx}
\caption{Selected major tunnelling projects in Australia. Engineers should verify current project status when citing external sources.}
\end{table}

\subsection{Infrastructure Risk Dashboard}

\paragraph{Critical Risks Identified}
The Infrastructure Risk Dashboard identifies geotechnical and ground movement issues as major risks for urban infrastructure delivery. For underground projects, these risks are often linked to uncertainty in geology, groundwater, existing structures, contaminated ground, land acquisition, and public acceptance.

\begin{table}[H]
\centering
\small
\begin{tabularx}{\textwidth}{X l l}
\toprule
\textbf{Sector risk} & \textbf{Typical rating} & \textbf{Typical trend} \\
\midrule
Ground movement and geotechnical issues & Critical & Increasing \\
Environmental and land contamination & High & Increasing \\
Social licence and community disruption & High & Increasing \\
Land acquisition and corridor constraints & Critical & Increasing \\
\bottomrule
\end{tabularx}
\caption{Infrastructure risk dashboard summary.}
\end{table}

\begin{keybox}[Engineering interpretation]
A risk rating is not only an administrative label. It affects investigation scope, contingency planning, construction methodology, monitoring requirements, and contractual allocation of risk.
\end{keybox}

\section{Australian project case studies}
\label{sec:australian-project-case-studies}

\subsection{Case Study 1: Melbourne Metro Tunnel}

\paragraph{Geological Setting}
The Melbourne Metro Tunnel is a useful example of a shallow urban underground project constructed in variable ground conditions. The dominant rock unit is the \textbf{Melbourne Formation}, comprising Silurian siltstone and sandstone, with weathering and strength varying significantly with depth and location.

\begin{table}[H]
\centering
\footnotesize
\begin{tabularx}{\textwidth}{l c c c c X c c}
\toprule
\textbf{Unit} & \boldmath{$\sigma_{ci}$} \textbf{(MPa)} & \boldmath{$m_i$} & \boldmath{$E_i$} \textbf{(GPa)} & \boldmath{$\nu$} & \textbf{Weathering} & \textbf{GSI} & \textbf{Depth (m)} \\
\midrule
MF1 & 19 & 13 & 6.0 & 0.20 & Slightly weathered--fresh & 55 & Below 30 \\
MF2 & 61 & 3 & 3.5 & 0.20 & Moderately weathered & 45 & 15--30 \\
MF3 & 21 & 3 & 1.5 & 0.25 & Highly weathered & 35 & 5--15 \\
MF4 & 1 & -- & 0.4 & 0.30 & Extremely weathered & -- & 0--5 \\
\bottomrule
\end{tabularx}
\caption{Representative Melbourne Metro geological units.}
\end{table}

\begin{notebox}[Important interpretation]
MF2 has a higher intact compressive strength than MF1 even though it is more weathered. This reflects lithological variation within the Melbourne Formation. Weathering grade and intact strength are not always inversely correlated when lithology changes.
\end{notebox}

\paragraph{Key Geotechnical Challenges}
For shallow urban tunnelling, the main concern is often not only tunnel collapse, but also the effect of excavation on the urban environment. Key challenges include:
\begin{enumerate}
 \item \textbf{Geological structures}: faults, folds, dykes, and variable bedding conditions;
 \item \textbf{Mixed ground conditions}: transitions between weathered rock, stronger rock, and soil-like material;
 \item \textbf{Existing infrastructure}: interaction with nearby buildings, basements, utilities, City Loop tunnels, and road tunnels;
 \item \textbf{Surface and groundwater}: potential influence of the Yarra River and shallow groundwater; and
 \item \textbf{Settlement control}: serviceability limits may govern design even when overall stability is adequate.
\end{enumerate}

\paragraph{Cost and Risk Context}

\begin{table}[H]
\centering
\small
\begin{tabularx}{0.8\textwidth}{X r}
\toprule
\textbf{Works package} & \textbf{Cost (\$ million)} \\
\midrule
Land acquisition & 781.2 \\
Early works & 365.3 \\
Tunnel and stations & 5,319.7 \\
Rail infrastructure & 1,154.7 \\
Rail systems & 1,384.5 \\
Other state costs & 2,023.5 \\
\midrule
\textbf{Total original estimate} & \textbf{11,028.9} \\
\bottomrule
\end{tabularx}
\caption{Melbourne Metro works package costs based on original budget estimate.}
\end{table}

\begin{keybox}[Link to geotechnical design]
In shallow urban tunnels, a small amount of ground movement may still be unacceptable if it affects rail assets, buried utilities, heritage buildings, or sensitive facilities. Therefore, geotechnical design must consider both \textbf{ultimate stability} and \textbf{serviceability performance}.
\end{keybox}

\subsection{Case Study 2: Snowy 2.0}

\paragraph{Project Overview}
Snowy 2.0 is a deep underground pumped-storage hydroelectric project. It provides a contrasting case to the Melbourne Metro Tunnel because the governing geotechnical issues are strongly influenced by depth, in-situ stress, groundwater pressure, and geological uncertainty.

\begin{table}[H]
\centering
\small
\begin{tabularx}{1\textwidth}{X l}
\toprule
\textbf{Parameter} & \textbf{Value} \\
\midrule
Installed capacity & 2,000 MW \\
Large-scale storage & 350,000 MWh \\
Waterway tunnel length & Approximately 27 km \\
Power station depth & Approximately 850 m underground \\
Original Snowy Scheme & 225 km of tunnels constructed between 1949 and 1972 \\
\bottomrule
\end{tabularx}
\caption{Key Snowy 2.0 project statistics.}
\end{table}

\paragraph{Key Geotechnical Challenges}
\begin{enumerate}
 \item \textbf{Geological uncertainty}: variable rock mass conditions along long tunnel alignments;
 \item \textbf{High in-situ stress}: potential stress-induced damage, spalling, or rockburst in deep excavations;
 \item \textbf{Groundwater}: complex hydrogeological conditions and potential high-pressure inflows;
 \item \textbf{Rock mass variability}: high-strength intact rock may still contain faults, shear zones, and weak zones; and
 \item \textbf{Construction access}: limited underground access increases the consequence of unexpected ground conditions.
\end{enumerate}

\begin{warnbox}[Snowy 2.0 as a geotechnical risk case]
Snowy 2.0 demonstrates how geological uncertainty can significantly influence cost, schedule, and construction strategy. In the design study, engineers should use this case to explain why early-stage investigation, parameter selection, and uncertainty management are central to underground engineering.
\end{warnbox}

\subsection{Case Study 3: West Gate Tunnel}

\begin{warnbox}[Environmental and contaminated ground risk]
The West Gate Tunnel illustrates that geotechnical risk is not limited to strength and deformation. Contaminated spoil, groundwater interaction, disposal requirements, and environmental approvals can dominate project cost and programme. These issues should be considered part of the broader geotechnical risk framework.
\end{warnbox}

\section{Shallow vs Deep Underground Conditions}

\begin{table}[H]
\centering
\small
\begin{tabularx}{\textwidth}{l X X}
\toprule
\textbf{Aspect} & \textbf{Melbourne Metro Tunnel: shallow urban tunnel} & \textbf{Snowy 2.0: deep underground project} \\
\midrule
Depth & Typically shallow to moderate & Up to several hundred metres below ground \\
Geology & Weathered Melbourne Formation, variable siltstone/sandstone & Granitic and metamorphic rocks with faults and weak zones \\
Weathering & Strongly depth-dependent & Generally fresher rock, but geological structures remain critical \\
Stress regime & Low to moderate & High; stress-induced failure may govern \\
Ground control & Settlement, mixed ground, water, interaction with assets & Rock mass instability, squeezing, spalling, groundwater pressure \\
Groundwater & Surface water and shallow groundwater interaction & Potentially high-pressure groundwater systems \\
Construction & TBM and cut-and-cover station works & TBM and drill-and-blast methods \\
Key design concern & Ground movement and serviceability & Stability, stress damage, support capacity, constructability \\
\bottomrule
\end{tabularx}
\caption{Comparison between shallow and deep underground conditions.}
\end{table}

\begin{keybox}[How to compare projects]
Do not simply list project facts. A good comparison explains how the differences in depth, stress, geology, groundwater, and urban constraints change the geotechnical design problem.
\end{keybox}

\section{Environmental approvals and decarbonisation}
\label{sec:environmental-approvals-and-decarbonisation}

\subsection{Environmental Assessment Framework}

\paragraph{EIS vs EES}
Large infrastructure projects in Australia require environmental assessment. The terminology varies between jurisdictions. Victoria commonly uses an \textbf{Environmental Effects Statement (EES)}, while other jurisdictions may use \textbf{Environmental Impact Statement (EIS)} terminology.

\begin{table}[H]
\centering
\small
\begin{tabularx}{\textwidth}{l X X}
\toprule
\textbf{Aspect} & \textbf{EIS} & \textbf{EES} \\
\midrule
Jurisdiction & Used in several Australian states and territories & Specific to Victoria \\
Legislation & State-specific environmental planning legislation & Environment Effects Act 1978 (Victoria) \\
Commonwealth trigger & \multicolumn{2}{X}{EPBC Act 1999 applies when Matters of National Environmental Significance may be affected.} \\
Typical relevance to geotechnics & Groundwater, spoil, contamination, settlement, vibration, ground movement & Groundwater, spoil, contamination, settlement, vibration, ground movement \\
\bottomrule
\end{tabularx}
\caption{Comparison between EIS and EES frameworks.}
\end{table}

\paragraph{EPBC Act and Key Tools}
The \textbf{Environment Protection and Biodiversity Conservation Act 1999} is the Australian Government's key environmental legislation for matters of national environmental significance.

\begin{itemize}
 \item The EPBC Act Public Portal can be used to search project referrals and approval status.
 \item The Protected Matters Search Tool can identify protected species, habitats, and environmental constraints.
 \item For underground projects, environmental assessment often includes groundwater, settlement, vibration, spoil disposal, contamination, and impacts on sensitive receptors.
\end{itemize}

\subsection{Net Zero, Energy Storage, and Underground Engineering}

Australia's net-zero transition has increased interest in renewable energy storage, including pumped hydro energy storage. Projects such as Snowy 2.0, Kidston, Borumba, and Pioneer-Burdekin require substantial underground engineering input.

\begin{keybox}[Energy transition link]
The energy transition creates new demand for underground engineering. Pumped hydro projects require tunnels, shafts, caverns, portals, and complex groundwater management. This creates direct links between climate policy and geotechnical engineering practice.
\end{keybox}

\section{Geotechnical Site Investigation and Rock Mass Classification}

\paragraph{Why Investigation Matters}
A rock mass is not the same as intact rock. It consists of intact blocks separated by discontinuities such as joints, bedding planes, faults, shears, and weathered zones. Its engineering behaviour is governed by both intact rock properties and discontinuity characteristics.

\begin{notebox}[Investigation-to-design chain]
A typical design chain is:
\[
\begin{aligned}
&\text{Field observations} \rightarrow \text{Core logging}
\rightarrow \text{RQD/RMR/Q/GSI} \\
&\rightarrow \text{Strength and deformation parameters}
\rightarrow \text{Design model}
\end{aligned}
\]
Each transformation introduces uncertainty. This idea becomes central in Chapters~\ref{ch:stress-intact-rock} to~\ref{ch:site-investigation}.
\end{notebox}

\paragraph{Anderson Creek Formation Example}
The Anderson Creek Formation siltstone and sandstone can be divided into subunits based on intact strength, degree of fracturing, and degree of weathering. Lower AF numbers generally indicate less weathered and better-quality rock.

\begin{table}[H]
\centering
\footnotesize
\begin{tabularx}{\textwidth}{l X X X X X X}
\toprule
 & \multicolumn{6}{c}{\textbf{Intact strength class}} \\
\cmidrule(lr){2-7}
\textbf{GSI group} & \textbf{Soil} & \textbf{VL} & \textbf{L} & \textbf{M} & \textbf{H} & \textbf{VH} \\
 & $<0.6$ MPa & $<0.6$ MPa & $2$--$6$ MPa & $6$--$20$ MPa & $20$--$60$ MPa & $>60$ MPa \\
\midrule
GSI1: 50--$>65$ & N/A & N/A & AF6 & AF4 & AF1 & AF1 \\
GSI2: 40--$>55$ & N/A & N/A & AF6 & AF4 & AF2 & AF2 \\
GSI3: 30--$>45$ & N/A & N/A & AF7 & AF5 & AF3 & N/A \\
GSI4: 20--$>35$ & N/A & N/A & AF7 & AF5 & AF3 & N/A \\
GSI5: 0--$>30$ & \multicolumn{6}{c}{Fault/shear zone} \\
\bottomrule
\end{tabularx}
\caption{Example rock mass class definition matrix for the Anderson Creek Formation.}
\end{table}
